\section*{Chapitre \ref{ch:b1:earthquake-source} --- La source sismique}

\begin{solution}{exo:b1:earthquake-source:1}
(a) $\log_{10}(\num{4e17}) = 17.60$, donc $M_w = \tfrac23(17.60 - 9.1) = 5.67$.
(b) $\log_{10}(\num{5.59e22}) = 22.75$ et $M_w = \tfrac23(22.75 - 9.1) =
9.10$. (c) $M_0 = 10^{1.5\times7.0 + 9.1} = 10^{19.6} = \qty{3.98e19}{N.m}$.
\end{solution}

\begin{solution}{exo:b1:earthquake-source:2}
(a) $T = D/\dot s = \qty{3}{m}/\qty{0.025}{m.\year^{-1}} = \qty{120}{\year}$.
(b) Avec $c = 0.5$, $T = D/(c\dot s) = \qty{240}{\year}$~: le fluage absorbe la moitié
du mouvement, et les séismes, de même glissement, viennent deux fois moins souvent.
\end{solution}

\begin{solution}{exo:b1:earthquake-source:3}
D'après la \cref{prop:b1:earthquake-source:moment-steps}, le rapport des moments est
$10^{1.5\times(7.4 - 6.2)} = 10^{1.8} = 63.1$~: il faut 63 séismes de
magnitude 6.2 pour libérer le moment d'un séisme de magnitude 7.4.
\end{solution}

\begin{solution}{exo:b1:earthquake-source:4}
Avec $b = 1$, chaque unité de magnitude divise le décompte par 10~: 10 séismes de
magnitude 5 ou plus par an, 1 de magnitude 6 ou plus, et 0.1 de magnitude 7 ou
plus, c'est-à-dire un tous les 10 ans en moyenne.
\end{solution}

\begin{solution}{exo:b1:earthquake-source:5}
(a) $A = \num{4.3e5} \times \qty{1.5e4}{m} = \qty{6.45e9}{m^2}$, et
$M_0 = \num{3e10}\times\num{6.45e9}\times4 = \qty{7.74e20}{N.m}$~;
$\log_{10}M_0 = 20.89$ et $M_w = \tfrac23(20.89 - 9.1) = 7.86$. (b) Une profondeur de
\qty{10}{km} divise l'aire et le moment par 1.5, ce qui abaisse la magnitude
de $\tfrac23\log_{10}1.5 = 0.12$, à 7.74. (c) La première, 7.86, s'arrondit au
7.9 du catalogue~; la seconde est un peu basse. La magnitude n'est pas très
sensible à la profondeur~: un facteur 1.5 sur l'aire ne fait que 0.12 de magnitude.
\end{solution}

\begin{solution}{exo:b1:earthquake-source:6}
(a) $D = M_0/(\mu A) = \num{5.59e22}/(\num{3e10}\times\num{4e5}\times\num{1.5e5})
= \qty{31.1}{m}$. (b) $\qty{31.1}{m}/\qty{0.083}{m.\year^{-1}} = \qty{374}{\year}$.
(c) Avec la fraction bloquée 0.8, $31.1/(0.8\times0.083) = \qty{468}{\year}$~; avec
0.5, \qty{749}{\year}. Moins l'interface est bloquée, plus le glissement
d'un grand séisme met de temps à se constituer~; le glissement moyen trouvé ici avec un
rectangle de \qtyproduct{400 x 150}{km} et $\mu = \qty{30}{GPa}$ n'est lui-même
qu'une estimation, et le vrai glissement était loin d'être uniforme.
\end{solution}

\begin{solution}{exo:b1:earthquake-source:7}
(a) La direction est $\phi = 329$, le \omterm{def:b1:geological-maps:attitude}{sens du pendage} $a = 059$, le \omterm{def:b1:geological-maps:attitude}{pendage} $\delta =
\ang{8}$. Alors, dans le repère est--nord--haut, $\mathbf s = (\sin329, \cos329, 0) = (-0.5150, 0.8572, 0)$,
$\mathbf u = (-\cos8\sin59, -\cos8\cos59, \sin8) = (-0.8488, -0.5100, 0.1392)$,
$\mathbf n = (\sin8\sin59, \sin8\cos59, \cos8) = (0.1193, 0.0717, 0.9903)$,
et avec $r = \ang{110}$,
\[
  \mathbf d = \cos110\,\mathbf s + \sin110\,\mathbf u = (-0.6215, -0.7724, 0.1308) .
\]
(b) Le \omterm{def:b1:earthquake-source:focal-mechanism}{plan auxiliaire} a pour normale $\mathbf d$ et pour glissement $\mathbf n$.
En prenant la normale vers le haut, il plonge de $\arccos 0.1308 = \ang{82.5}$ vers
$\arctan(-0.6215/-0.7724) = 218.8$, donc sa direction est $128.8$~; le glissement
$\mathbf n$ projeté sur ses vecteurs de direction et de remontée du \omterm{def:b1:geological-maps:attitude}{pendage} donne le \omterm{def:b1:faults-folds-fabrics:fault-elements}{pitch}
\ang{87.2}. Arrondi~: $129/83/87$, le second plan de la solution publiée.
(c) $\mathbf P = (\mathbf n - \mathbf d)/\sqrt2$ a pour direction $221$ et plonge de
\ang{37}~; $\mathbf T = (\mathbf n + \mathbf d)/\sqrt2$ a pour direction $036$ et plonge de
\ang{52} (publié~: $222/38$ et $36/52$). (d) Le plan peu profond, plongeant de
\ang{8} vers $059$, le long du \omterm{def:b1:active-margins:slab}{panneau plongeant}~: une \omterm{def:b1:faults-folds-fabrics:fault-types}{faille inverse} de \omterm{def:b1:faults-folds-fabrics:fault-elements}{pitch}
\ang{110}, et, plongeant de moins de \ang{45}, un \omterm{def:b1:faults-folds-fabrics:fault-types}{chevauchement}, l'interface de plaques de
la zone de subduction.
\end{solution}

\begin{solution}{exo:b1:earthquake-source:8}
(a) Le plan plonge vers $135$, donc $\mathbf n = (\sin135, \cos135, 0) =
(0.707, -0.707, 0)$~; un glissement \omterm{def:b1:faults-folds-fabrics:fault-types}{dextre} a $r = \ang{180}$, $\mathbf d = -\mathbf s
= (-0.707, -0.707, 0)$. Vers le nord, $\mathbf x = (0, 1, 0)$~:
$(\mathbf n\cdot\mathbf x)(\mathbf d\cdot\mathbf x) = (-0.707)(-0.707) > 0$, une
poussée. Est~: $(0.707)(-0.707) < 0$, une traction. Sud~: une poussée~; ouest~: une traction.
(b) $\mathbf P = (\mathbf n - \mathbf d)/\sqrt2 = (1, 0, 0)$~: l'\omterm{def:b1:earthquake-source:focal-mechanism}{axe P} est
horizontal de direction $090$--$270$~; $\mathbf T = (0, -1, 0)$ est de direction
$000$--$180$. (c) Le plan vertical de direction $135$--$315$ a pour normale
$\mathbf d$. Avec le glissement $\mathbf n$, le bloc de son côté sud-ouest, vers
lequel pointe $\mathbf d$, se déplace vers $135$~; vu du bloc nord-est,
en regardant vers le sud-ouest, c'est vers la gauche~: un glissement \omterm{def:b1:faults-folds-fabrics:fault-types}{senestre}. D'après la
\cref{prop:b1:earthquake-source:auxiliary}, ce glissement donne les mêmes \omterm{def:b1:earthquake-source:focal-mechanism}{premiers
mouvements} dans toutes les directions.
\end{solution}

\begin{solution}{exo:b1:earthquake-source:9}
(a) $V = \qty{40}{mm.\year^{-1}}$ et $v = V/4$ à \qty{10}{km}, donc d'après la
\cref{prop:b1:earthquake-source:profile}(b) $D = \qty{10}{km}$~; directement,
$D = x/\tan(\pi v/V) = 10/\tan(\pi/4) = \qty{10}{km}$. (b) $v(30) =
(40/\pi)\arctan 3 = \qty{15.9}{mm.\year^{-1}}$. (c) $V/(\pi D) =
\num{0.040}/(\pi\times\num{10000}) = \num{1.27e-6}$ par an.
\end{solution}

\begin{solution}{exo:b1:earthquake-source:10}
(a) Les moments des classes sont toujours $C q^k$ avec $q = 10^{1.5 - b}$, mais maintenant
$q < 1$~: le moment des classes décroît vers les grandes magnitudes, et la
classe inférieure, $k = 0$, porte la part $C/S = (1 - q)/(1 - q^{K+1}) > 1 - q$,
image en miroir de la \cref{prop:b1:earthquake-source:largest}. Les plus petits
séismes comptés domineraient alors, et le total croîtrait sans limite
à mesure que l'on inclurait des séismes de plus en plus petits. (b) $1 - 10^{-0.7} = 0.80$~: au moins
\qty{80}{\percent}. (c) $\log_{10}N = 8.20 - 1.01\times9 = -0.89$, donc $N = 0.13$
par an et environ 6.4 en cinquante ans, trois fois les deux observés. La droite,
ajustée entre les magnitudes 6 et 7.5, surestime les plus grands séismes~;
ou bien la loi s'infléchit vers le bas au sommet, comme le suggère la taille finie des failles, ou
bien cinquante ans sont trop peu pour échantillonner des événements qui viennent une fois par décennies. Les deux sont
possibles avec deux événements.
\end{solution}

\begin{solution}{exo:b1:earthquake-source:11}
(a) En 150 ans un poteau en $x$ se déplace de $150\,v(x) = 150\times(34/\pi)
\arctan(x/15)$ millimètres par rapport au poteau sur la faille~: \qty{0.11}{m} à
\qty{1}{km}, \qty{1.28}{m} à \qty{15}{km}, \qty{2.08}{m} à \qty{50}{km},
vers le sud pour une faille \omterm{def:b1:faults-folds-fabrics:fault-types}{dextre} vue de l'ouest. (b) Le déficit de glissement est
$s = \qty{0.034}{m.\year^{-1}}\times\qty{150}{\year} = \qty{5.1}{m}$. Pendant le
séisme les poteaux se déplacent de $(s/\pi)\arctan(D/x)$~: \qty{2.44}{m} à
\qty{1}{km}, \qty{1.28}{m} à \qty{15}{km}, \qty{0.47}{m} à \qty{50}{km}.
(c) Les sommes sont $0.11 + 2.44$, $1.28 + 1.28$ et $2.08 + 0.47$, toutes
\qty{2.55}{m} $= s/2$~: après le séisme la clôture est de nouveau droite du côté
est, déplacée de $s/2$, et le côté ouest, par symétrie, de $-s/2$.
\end{solution}

\begin{solution}{exo:b1:earthquake-source:12}
(a) $\mathbf P = (0, 1, 0)$ et $\mathbf T = (1, 0, 0)$. Le couple
$\mathbf n = (\mathbf T + \mathbf P)/\sqrt2 = (0.707, 0.707, 0)$,
$\mathbf d = (\mathbf T - \mathbf P)/\sqrt2 = (0.707, -0.707, 0)$ vérifie
$\mathbf P = (\mathbf n - \mathbf d)/\sqrt2$ et $\mathbf T = (\mathbf n +
\mathbf d)/\sqrt2$. Les deux vecteurs sont horizontaux, donc les deux \omterm{def:b1:earthquake-source:focal-mechanism}{plans nodaux} sont
verticaux~: le plan de normale $\mathbf n$ (vers $045$) a pour direction $135$--$315$,
le plan de normale $\mathbf d$ (vers $135$) a pour direction $045$--$225$. (b) Sur le
plan de direction $315$, le bloc du côté nord-est se déplace selon $\mathbf d$,
vers $135$~: vu de l'autre bloc, vers la droite, un glissement \omterm{def:b1:faults-folds-fabrics:fault-types}{dextre}. Sur le
plan de direction $045$, le bloc du côté sud-est se déplace selon $\mathbf n$,
vers $045$~: un glissement \omterm{def:b1:faults-folds-fabrics:fault-types}{senestre}. (c) Une nouvelle faille fait $\ang{45} - \phi/2 =
\ang{30}$ avec $\sigma_1$~: avec $\sigma_1$ de direction $000$, les failles auraient pour
directions $030$ (\omterm{def:b1:faults-folds-fabrics:fault-types}{senestre}) et $330$ (\omterm{def:b1:faults-folds-fabrics:fault-types}{dextre}), à \ang{15} des \omterm{def:b1:earthquake-source:focal-mechanism}{plans nodaux}.
L'\omterm{def:b1:earthquake-source:focal-mechanism}{axe P} est fixé à \ang{45} des \omterm{def:b1:earthquake-source:focal-mechanism}{plans nodaux} par la géométrie, quel que soit le
frottement de la roche, donc il coïncide avec $\sigma_1$ seulement si la faille a glissé
à \ang{45} de lui. (d) La faille de direction $040$ est à \ang{5} du \omterm{def:b1:earthquake-source:focal-mechanism}{plan nodal}
de direction $045$~: ce plan est vraisemblablement la faille, et le glissement sur lui est
\omterm{def:b1:faults-folds-fabrics:fault-types}{senestre}.
\end{solution}

\begin{solution}{pb:b1:earthquake-source:1}
\textbf{1.} Les intervalles sont $1290 - 1050 = 240$, $1480 - 1290 = 190$ et
$1712 - 1480 = 232$ ans~; leur moyenne est \qty{220.7}{\year}.

\textbf{2.} Le décalage moyen est $(7.0 + 6.0 + 5.5 + 7.5)/4 = \qty{6.5}{m}$, et
$T = 6.5/0.030 = \qty{216.7}{\year}$, proche de l'intervalle moyen.

\textbf{3.} Les décalages des trois derniers séismes, $6.0 + 5.5 + 7.5 =
\qty{19.0}{m}$, se sont accumulés entre 1050 et 1712, en 662 ans~: $\qty{19.0}{m}
/\qty{662}{\year} = \qty{28.7}{mm.\year^{-1}}$, \qty{4}{\percent} de moins que les
\qty{30}{mm.an^{-1}} géodésiques. Les deux concordent dans l'incertitude des
dates et des décalages~: le segment libère presque tout son chargement dans ses grands
séismes.

\textbf{4.} $2025 - 1712 = 313$ ans, et le déficit est
$\qty{0.030}{m.\year^{-1}}\times\qty{313}{\year} = \qty{9.39}{m}$.

\textbf{5.} Le déficit dépasse tous les décalages de la tranchée, le plus grand étant
\qty{7.5}{m}, et le temps depuis 1712 dépasse le plus long intervalle, 240 ans.
En ce sens le segment est en retard~; cela ne veut pas dire que le séisme est
imminent, seulement que la faille détient déjà plus de glissement qu'elle n'en a jamais
libéré d'un coup dans l'enregistrement.

\textbf{6.} Les intervalles vont de 190 à 240 ans et les décalages de
\qtyrange{5.5}{7.5}{m}~: le cycle n'est pas périodique, et la résistance de la
faille au prochain séisme est inconnue. La tranchée donne une moyenne et une dispersion,
d'où l'on peut estimer une probabilité, non une date.

\textbf{7.} $A = \qty{250}{km}\times\qty{15}{km} = \qty{3750}{km^2} =
\qty{3.75e9}{m^2}$.

\textbf{8.} $M_0 = \num{3e10}\times\num{3.75e9}\times\num{9.39} =
\qty{1.06e21}{N.m}$.

\textbf{9.} $M_0 = \num{3e10}\times\num{3.75e9}\times\num{7.5} =
\qty{8.44e20}{N.m}$, so $M_w = \tfrac23(20.926 - 9.1) = 7.88$.

\textbf{10.} Le rapport est $9.39/7.5 = 1.25$, ce qui fait $\tfrac23\log_{10}1.25 =
0.07$ de magnitude.

\textbf{11.} La moitié de l'aire divise le moment par deux et abaisse la magnitude de
$\tfrac23\log_{10}2 = 0.20$~: $M_w = 7.75$.

\textbf{12.} $\mu A V = \num{3e10}\times\num{3.75e9}\times\num{0.030} =
\qty{3.38e18}{N.m}$ par an.

\textbf{13.} $b = 1$. $N(5) = 10^{4.3 - 5} = 0.20$ par an et $N(6) = 0.020$ par
an.

\textbf{14.} $N(7) = 0.0020$ par an~: un tous les 500 ans~; $N(8) = 0.00020$ par
an~: un tous les \num{5000} ans.

\textbf{15.} La loi donne un séisme de magnitude 7.95 ou plus tous les
$1/10^{4.3 - 7.95} = \num{4500}$ ans, la tranchée un tous les 220 ans, vingt
fois plus souvent. La loi est ajustée aux nombreux petits séismes de toutes les
failles de la région~; extrapolée aux grands séismes d'un seul segment, elle
les sous-estime fortement. La récurrence des grands séismes sur une faille
se tire de son propre enregistrement et de son chargement, non des petits.

\textbf{16.} $\num{5.5e17}/\num{3.38e18} = 0.16$~: les séismes modérés
libèrent \qty{16}{\percent} du moment que le segment accumule.

\textbf{17.} D'après la \cref{prop:b1:earthquake-source:moment-steps},
$10^{1.5\times(7.95 - 6)} = 10^{2.92} \approx 840$ fois plus.

\textbf{18.} Avec $b = 1$, le moment des séismes est dominé par les
plus grands~: quoi que fassent les séismes modérés, ils libèrent une part mineure, ici
\qty{16}{\percent}, et les \qty{84}{\percent} restants, le déficit de glissement,
attendent le grand séisme.

\textbf{19.} $v(-x) = -v(x)$ pour chaque paire, donc le profil est antisymétrique.
Loin de la faille les vitesses tendent vers $\pm V/2 = \pm\qty{15}{mm.\year^{-1}}$.

\textbf{20.} $D = x/\tan(\pi v/V) = 15/\tan(\pi\times7.5/30) = 15/\tan\ang{45} =
\qty{15}{km}$~: à la \omterm{def:b1:earthquake-source:locking-depth}{profondeur de blocage} la vitesse est $V/4$.

\textbf{21.} À \qty{5}{km}~: $D = 5/\tan(\pi\times3.1/30) = 5/\tan\ang{18.6} =
\qty{14.9}{km}$~; à \qty{40}{km}~: $D = 40/\tan\ang{69.6} = \qty{14.9}{km}$. Les
vitesses sont données à \qty{0.1}{mm.an^{-1}}, et cet arrondi seul déplace
$D$ d'environ \qty{0.1}{km}~; avec des données réelles, le bruit de chaque station serait
plus grand, et la \omterm{def:b1:earthquake-source:locking-depth}{profondeur de blocage} serait ajustée à toutes les stations à la fois.

\textbf{22.} La \omterm{def:b1:earthquake-source:locking-depth}{profondeur de blocage}, \qty{15}{km}, égale la profondeur des plus profonds
séismes. La partie bloquée de la faille est sa partie cassante, la \omterm{def:b1:stress-strain-rheology:strength-envelope}{couche
sismogène}, où le frottement la retient entre les séismes et cède dans eux~; en dessous, la
roche plus chaude flue (\cref{ch:b1:stress-strain-rheology}) et la faille glisse
régulièrement, ce qui est le glissement à la vitesse $V$ sous $D$ du modèle.

\textbf{23.} $V/(\pi D) = \num{0.030}/(\pi\times\num{15000}) = \num{6.4e-7}$
par an.

\textbf{24.} Avec $M_0 = \qty{1.056e21}{N.m}$, $\log_{10}M_0 = 21.024$ et
$M_w = \tfrac23(21.024 - 9.1) = 7.949$. À 0.01 près, le séisme qui libérerait
tout le déficit du segment a la \omterm{def:b1:earthquake-source:mw}{magnitude de moment}
\textbf{7.95}.
\end{solution}
